\section{Conclusions}

The ALICE Collaboration has conducted comprehensive studies of the production of strange hadrons in various hadronic collision systems, 
revealing a remarkable smooth trend that correlates with increasing multiplicity from pp to heavy-ion collisions. Ongoing analyses and investigations 
to understand the origins of strangeness production and strangeness enhancement in hard and soft processes are of particular interest, and the results 
presented here shed some light on these particle production mechanisms in small collision systems.

Utilising the method of two particle correlations, it can be observed that strange hadrons observed in jets have different production mechanism than 
those outside the jet cone, based on the harder transverse momentum spectra for in-jet particles. Initial collision energy dependence has not been 
observed for strangeness enhancement in any region, however dependence on charged particle multiplicity is noticed. The per-trigger yield 
ratios demonstrate this multiplicity dependent enhancement, which is compatible in jet and out-of jet regions. However, all observed cases point to 
a larger contribution from out-of-jet regions in comparison to in-jet region, indicating soft processes dominate over hard scattering processes 
in strange particle production.

Further studies based on dataset from Run 3 of the LHC are anticipated to provide more extensive and detailed outcomes thanks to a significantly larger 
sample of high-multiplicity pp collisions.
